% Authoritative LaTeX manuscript; Markdown is a generated compatibility projection.

% <!-- part-title: The New Engineering Problem -->

% <!-- chapter-title: The Printer -->

\chapter{The Printer}

% <!-- point: agents-made-implementation-cheap-not-reliable | Coding agents make software implementation cheap enough to change the engineering problem, creating the question of what engineers should do with the new productive capacity. | terms: code-fabricator -->

Coding agents make software implementation cheap enough to change the engineering problem. Describe a change and an agent can often implement it quickly. The opportunity is substantial, but what should we do with this new capacity?

% <!-- index-def: printer-metaphor -->

For an analogy, let us begin with office supplies. Consider first a stapler and a hole punch. Their purposes and limitations are in their names. Need two sheets joined together, or a hole punched? Look no further. These tools perform narrow operations on artifacts that already exist.

% <!-- point: the-3d-printer-is-open-ended | A printer can produce many kinds of artifacts, subject to the capabilities and limitations of its machinery. | terms: printer-metaphor -->

% <!-- point: open-ended-capacity-must-be-engineered | People make broad productive capacity useful by developing forms and ways of working. | terms: printer-metaphor, commodity-intelligence -->

Printers are different. A 2D printer can produce pages of text and images. A 3D printer can build physical objects from descriptions of their shapes. The same machine can produce many different artifacts, subject to the capabilities and limitations of the machinery. People make that productive capacity useful by developing forms and ways of working. For 2D printers, familiar forms include newspapers, posters, and maps, each with conventions for layout and legibility. For a 3D printer, you might be making a mounting bracket, a protective enclosure, or an adapter. These have different purposes and recurring requirements for strength, clearance, and fit. Templates, reusable designs, and established production practices give people starting points for making such artifacts. People develop and adapt this knowledge as they discover what the machinery can usefully make.

% <!-- point: agent-inference-needs-engineered-boundaries | Software agents supply broad productive capacity and can infer implementation decisions that incomplete instructions leave unspecified. | terms: printer-metaphor, commodity-intelligence -->

Software agents are printers, not staplers. They likewise supply broad productive capacity. Like printers, agents can follow instructions. But agents can infer much more from an incomplete description than a physical printer can: an algorithm, a data structure, or a decomposition may never have been specified, and yet an agent can choose them and produce an implementation. That inferential capacity is precisely what makes the machinery powerful.

% <!-- point: delegating-judgment-is-an-opportunity-not-a-requirement | Delegation can leave useful judgment to the agent, but the engineer decides which decisions to make explicit and which to delegate. -->

When we work with agents, we delegate some judgment as well as production. We could specify the entire implementation, hand it to an agent, and ask it to reproduce it. That would be an unusually expensive way to copy a file. Much of the opportunity lies in letting agents work out things we have not already worked out ourselves. The engineering question is which decisions we should make explicit and which we can responsibly leave to the agent.
